\section{Aligning the modular action}
\label{sec:modular-alignment}

The preceding section compares the spectral subspaces of two modular
operators. We now turn that comparison into an algebra invariant under
one modular group. Distances are measured in operator norm, and all
constants are absolute. The selected unitaries need not depend
continuously on time; the proof specifies the topology whenever
continuity is used.

\begin{proposition}[A common modular action]
\label{prop:modular-alignment}
There are absolute constants $c_{\rm mod},C_{\rm mod}>0$ with the
following property. Let $(M,N,\Omega,R)$ be a prepared pair as in
\Cref{def:prepared-pair}, acting on the same separable Hilbert space
$H$, and suppose that $\dcb(M,N)<\gamma<c_{\rm mod}$.
There is a unitary $q\in\B(H)$ such that
\[
 \norm{q-I}\le C_{\rm mod}\gamma,
 \qquad \dcb(qMq^*,N)\le C_{\rm mod}\gamma,
 \qquad
 \Delta_N^{it}(qMq^*)\Delta_N^{-it}=qMq^*
 \quad(t\in\R).
\]
The algebra $N$ and its modular group are unchanged.
\end{proposition}

The proof has three steps. A uniform infinitesimal estimate controls
norm-continuous paths of approximate normalizers. Analytic smoothing
puts the modular comparison into that setting, and a completely
positive limit recovers exact normalizers at each time. Finally an
averaging iteration corrects their multiplication law and produces
the near-identity change of position.

We use two groups associated with a concrete factor $P\subset\B(H)$:
\[
 \mathcal G_P=\U(P)\U(P'),\qquad
 \mathcal N_P=\{v\in\U(\B(H)):vPv^*=P\}.
\]
The two factors in a product defining $\mathcal G_P$ commute, so
$\mathcal G_P$ is a subgroup of $\mathcal N_P$. We first prove a
uniform estimate near this subgroup. This estimate uses only that $P$
is a separably acting type $\III$ factor.

\subsection{Distance to the infinitesimal normalizers}

\begin{lemma}
\label{lem:infinite-commutant-bound}
For a separably acting properly infinite von Neumann algebra
$P\subset\B(H)$ and $h=h^*\in\B(H)$,
\[
 \dist(h,P')\le4\norm{\ad_h|_P}.
\]
\end{lemma}

\begin{proof}
Put $d=\norm{\ad_h|_P}$. Choose a unital copy $S$ of
$\B(\ell^2)$ in $P$, with matrix units $(e_{ij})$ whose diagonal
sum is $I$. The finite-dimensional unital algebras
\[
 S_n=\operatorname{span}\{e_{ij}:1\le i,j\le n\}
       +\C\Bigl(I-\sum_{i=1}^ne_{ii}\Bigr)
\]
increase and generate $S$. Average $h$ over $\U(S_n)$ and take a
weak-operator cluster point $h_0$. Each average is self-adjoint,
within $d$ of $h$, and has scalar commutator bound at most $d$ on
$P$. These properties pass to the cluster point. Haar invariance
at every earlier stage gives $h_0\in S'$.

For any $k$, choose isometries $s_1,\ldots,s_k\in S$ with
$s_i^*s_j=\delta_{ij}I$ and $\sum_i s_is_i^*=I$.
The unitary $V:H^k\to H$, $V(\xi_i)=\sum_i s_i\xi_i$, identifies
$M_k(P)$ with $P$ and $h_0^{(k)}$ with $h_0$. Consequently
\[
 \norm{\ad_{h_0}|_P}_{\cb}\le d.
\]
By \Cref{lem:commutant-cb}, $\dist(h_0,P')\le d$, so in fact
$\dist(h,P')\le2d$. The stated larger constant will be convenient.
\end{proof}

\begin{lemma}[Uniform infinitesimal estimate]
\label{lem:local-normalizer-linear}
There is an absolute $K<\infty$ such that every separably acting
type $\III$ factor $P\subset\B(H)$ and every self-adjoint
$h\in\B(H)$ satisfy
\begin{equation}
 \dist(h,P+P')
 \le K\sup_{x\in P_1}\dist([h,x],P).
 \label{eq:normalizer-linear-bound}
\end{equation}
The sum $P+P'$ here is a linear subspace; no closedness assertion
about that sum is needed.
\end{lemma}

\begin{proof}
Suppose there is no common $K$. By scaling a counterexample and
subtracting a self-adjoint sum from $P+P'$, we can choose factors
$P_j\subset\B(H_j)$ and self-adjoint $h_j$ with
\[
 \dist(h_j,P_j+P_j')=1,
 \qquad \norm{h_j}\le2,
 \qquad
 \varepsilon_j:=\sup_{x\in(P_j)_1}\dist([h_j,x],P_j)\longrightarrow0.
\]
Taking real parts justifies using a self-adjoint approximating sum.
Its subtraction does not change either distance in the assertion,
because its commutator with $P_j$ belongs to $P_j$.

Fix a free ultrafilter $\omega$ on $\N$ and form the norm
ultraproducts
\[
 A=\prod_\omega P_j\subset \prod_\omega\B(H_j)=E.
\]
Here a norm ultraproduct is the algebra of bounded sequences modulo
the sequences whose operator norms tend to zero along $\omega$.
The algebra $A$ is simple. Indeed, a nonzero positive element $a$
has uniformly bounded positive representatives $a_j$ and a number
$t>0$ such that $e_j=1_{[t,\infty)}(a_j)$ is nonzero on a set in
$\omega$. In a separably acting type $\III$ factor each nonzero
projection is equivalent to $I$. Choose isometries $v_j$ with
$v_jv_j^*=e_j$ on that set. Then
$v_j^*a_jv_j\ge tI$ there. The class $v^*av$ is invertible in $A$
and belongs to the ideal generated by $a$. Thus every nonzero
closed ideal of $A$ contains $I$.

Commutation with $[h_j]\in E$ defines a bounded derivation
$D:A\to A$. To see that its range is in $A$, choose, for every
bounded sequence $x_j\in P_j$, elements $y_j\in P_j$ with
\[
 \norm{[h_j,x_j]-y_j}
 \le (\varepsilon_j+j^{-1})\sup_i\norm{x_i}.
\]
The $y_j$ are uniformly bounded, and their class is the commutator
in $E$. Its definition is therefore independent of the choices;
linearity and the derivation identity follow in $E$.

Sakai's theorem for simple unital $C^*$-algebras gives $a\in A$
with $D(x)=[a,x]$ \cite[Section~1, p.~259; Section~2, p.~260]{Sakai1971}.
No separability of the norm ultraproduct is required. Since the
$h_j$ are self-adjoint, replacing $a$ by $(a+a^*)/2$ leaves this
identity unchanged. Choose uniformly bounded self-adjoint
representatives $a_j\in P_j$. The equality of the two derivations
implies
\[
 \lim_\omega\norm{\ad_{h_j-a_j}|_{P_j}}=0.
\]
Otherwise contractions witnessing a fixed positive commutator norm
could be selected at each coordinate of an ultrafilter-large set,
contradicting equality on their class in $A$.
\Cref{lem:infinite-commutant-bound} now gives
\[
 1=\dist(h_j,P_j+P_j')
 \le\dist(h_j-a_j,P_j')
 \le4\norm{\ad_{h_j-a_j}|_{P_j}}\longrightarrow_\omega0,
\]
a contradiction.
\end{proof}

\begin{lemma}[A path of approximate normalizers]
\label{lem:local-normalizer-path}
There are absolute constants $e_0,C_{\rm path}>0$ such that the
following holds. Let $P\subset\B(H)$ be a separably acting type
$\III$ factor,
let $I\subset\R$ be an interval containing $0$, and let
$v:I\to\U(\B(H))$ be norm continuous with $v(0)=I$.
If $0\le e<e_0$ and
\[
 \sup_{t\in I}\sup_{x\in P_1}\dist(v(t)xv(t)^*,P)\le e,
\]
then
\[
 \sup_{t\in I}\dist(v(t),\mathcal G_P)\le C_{\rm path}e.
\]
The constants are independent of the interval's length.
\end{lemma}

\begin{proof}
Write $b(t)=\dist(v(t),\mathcal G_P)$. This is a continuous
function, since distance to any nonempty subset is $1$-Lipschitz.
Suppose $b=b(t)$ is smaller than a fixed absolute radius, and
choose $p\in\mathcal G_P$ with $\norm{p^*v(t)-I}<b+\eta$.
For sufficiently small $b+\eta$, principal functional calculus
writes $p^*v(t)=e^{ih}$ with $h=h^*$ and
$\norm{h}\le2(b+\eta)$. As $p$ normalizes $P$, this conjugation
has the same one-sided approximation bound $e$.
The twice-integrated conjugation formula gives
\[
 \norm{e^{ih}xe^{-ih}-x-i[h,x]}\le2\norm{h}^2\norm{x}.
\]
It follows that
\[
 \sup_{x\in P_1}\dist([h,x],P)\le e+2\norm{h}^2.
\]
Apply \Cref{lem:local-normalizer-linear} and approximate $h$ by
a self-adjoint sum $a+b'$ with $a\in P$, $b'\in P'$.
Since $a$ and $b'$ commute, $e^{i(a+b')}=e^{ia}e^{ib'}$ belongs
to $\mathcal G_P$. The inequality
$\norm{e^{ih}-e^{ik}}\le\norm{h-k}$ for self-adjoint $h,k$
therefore shows, on letting both approximation slacks vanish, that
\begin{equation}
 b(t)\le K e+8K b(t)^2
 \quad\text{whenever }b(t)\text{ is in that radius}.
 \label{eq:normalizer-path-barrier}
\end{equation}
Choose $r>0$ within the radius and with $8Kr\le1/4$.
If $Ke<r/2$, the continuous function $b$, starting at zero,
cannot first reach $r$: \eqref{eq:normalizer-path-barrier} would
give $r<3r/4$. Thus $b(t)<r$ throughout the interval between
$0$ and any chosen $t\in I$. Absorbing the quadratic term in
\eqref{eq:normalizer-path-barrier} gives $b(t)\le2Ke$ uniformly.
This proves the assertion with suitably reduced $e_0$.
\end{proof}

\subsection{Analytic approximation of the modular comparison}

The next argument converts separated spectral-cut estimates into actual
normalizers.  Its analytic approximation takes place in operator norm away
from the real axis.  Its passage to the real axis uses the strong-* topology
at each fixed time; it does not require norm continuity of modular unitaries.

We first record explicitly the analytic fact needed below.  On
$\mathcal K=L^2(\R,H)$ put $D=-i\,d/dt$, with Fourier-transform convention
\[
 \widehat f(\xi)=(2\pi)^{-1/2}\int_\R e^{-it\xi}f(t)\,dt,
 \qquad P_a=1_{(a,\infty)}(D).
\]
A bounded operator-valued function always means a field measurable on each
vector, modulo a common null set, with essentially bounded operator norm.
Here and below $H$ is separable.  Integrals of such fields are defined by
matrix coefficients, equivalently by integration on each vector when the
kernel is integrable.  No norm measurability as a $\B(H)$-valued function is
required.

\begin{lemma}[Analytic extension of a frequency-preserving multiplier]
\label{lem:analytic-multiplier}
Suppose multiplication by a bounded field $B(t)$ on $L^2(\R,H)$ preserves
$P_0\mathcal K$.  Then it is the boundary field of a bounded, operator-norm
holomorphic function on the upper half-plane, given by
\begin{equation}\label{eq:analytic-poisson}
 B(t+ir)=\int_\R \frac{r}{\pi((t-s)^2+r^2)}B(s)\,ds,
 \qquad r>0.
\end{equation}
Its norm is at most $\mathop{\rm ess\,sup}_s\norm{B(s)}$.
For two such fields, and any constant $x\in\B(H)$, the extension of the
boundary product $B(s)xS(s)$ is $B(z)xS(z)$.
\end{lemma}

\begin{proof}
Write $w=(t-i)/(t+i)$ and give the unit circle normalized measure
$d\theta/(2\pi)$.  The transformation
\[
 (Jf)(t)=\frac{1}{\sqrt\pi(t+i)}
             f\left(\frac{t-i}{t+i}\right)
\]
is unitary from $L^2(\mathbb T,H)$ onto $\mathcal K$: the change of variables
has $d\theta/(2\pi)=dt/(\pi(1+t^2))$.
For $n\geq0$, $J(w^n)$ is a linear combination of $(t+i)^{-k}$,
$1\leq k\leq n+1$, whose Fourier transforms are supported in $[0,\infty)$.
For $n<0$ its expression is a linear combination of $(t-i)^{-k}$, whose
Fourier transforms are supported in $(-\infty,0]$.  These support statements follow directly from
\[
 (t+i)^{-k}=\frac{(-i)^k}{(k-1)!}
       \int_0^\infty u^{k-1}e^{-u}e^{itu}\,du
\]
and its complex conjugate, interpreted under the $L^2$ Fourier transform.  Since the images of all circle modes form an orthonormal basis,
$J$ identifies the nonnegative-mode subspace with $P_0\mathcal K$.

Conjugating the multiplier by $J$ gives the circle multiplier
$\widetilde B(w)=B(t)$.  Apply its invariant-subspace property to every
constant vector.  All negative Fourier coefficients of $\widetilde B$ vanish.
If $b_n$ denotes its $n$th coefficient, defined by matrix coefficients, then
\[
 F(w)=\sum_{n\geq0} b_nw^n,\qquad |w|<1,
\]
converges in operator norm locally uniformly and is operator-norm
holomorphic.  Its value at $w=\rho e^{i\theta}$ is the circle Poisson
integral of $\widetilde B$, since all its negative coefficients vanish.
The Poisson kernel is positive and has integral one, so
$\norm{F(w)}\leq\norm{\widetilde B}_\infty$.
Changing variables in the circle Poisson integral gives exactly
\eqref{eq:analytic-poisson}; explicitly, harmonic measure at $t+ir$
pushes forward to $r\,ds/(\pi((t-s)^2+r^2))$.

For completeness, multiplication of these extensions can be checked
without presupposing pointwise boundary limits.  The radial Poisson
integrals converge in $L^2(\mathbb T,H)$ on each fixed vector, for both
$\widetilde B$ and $\widetilde B^*$, and likewise for $\widetilde S$ and
$\widetilde S^*$.  Choose a sequence of radii tending to one along which
these convergences hold pointwise for a countable dense family of vectors.
Uniform boundedness extends the pointwise convergence to every vector.
Along this sequence the products converge strongly-* almost everywhere
to $\widetilde B(w)x\widetilde S(w)$.  Each Fourier coefficient of the
bounded analytic product $F(w)xG(w)$ therefore converges, by dominated
convergence on matrix coefficients, to the corresponding coefficient of
that boundary product.  In particular its negative coefficients vanish,
and its nonnegative coefficients are precisely the Taylor coefficients
of $F xG$.  This proves the asserted product identity.  The same argument
applies to products with a scalar bounded analytic function.
\end{proof}

\begin{lemma}[From modular cuts to near normalizers]
\label{lem:modular-near-normalizers}
Let $M,N\subseteq\B(H)$ be a prepared pair on a separable Hilbert space,
with $M$ of type $\III$ and $\dcb(M,N)\leq\gamma$.  Write
\[
 U_P(t)=e^{itL_P},\qquad L_P=\log\Delta_P\quad(P=M,N),
 \qquad W(t)=U_N(t)U_M(t)^*.
\]
Assume the separated-cut estimate of \Cref{cor:modular-cuts}, with its
universal constant, and the local path estimate of
\Cref{lem:local-normalizer-path}.  There are universal constants
$\gamma_0>0$ and $C$ such that, if $0<\gamma<\gamma_0$, then for every
$t\in\R$ there is a unitary $n_t$ satisfying
\[
 n_tMn_t^*=M,\qquad \norm{n_t-W(t)}\leq C\gamma.
\]
No regularity of the choices $t\mapsto n_t$ is asserted or needed.
\end{lemma}

\begin{proof}
The letters $C$ below denote constants depending only on the universal
constants in the stated inputs, and can increase between occurrences.
All required smallness conditions are met by decreasing one universal
$\gamma_0$ at the end.

\paragraph{The initial approximate normalizers.}
Both $U_M(t)$ and $U_N(t)$ normalize their respective algebras.
For every matrix size $k$ and $x\in M_k(M)_1$, first move
$U_M(t)^*xU_M(t)$ into $M_k(N)_1$, conjugate by $U_N(t)$, and move back
into $M_k(M)_1$.  The two approximation errors are at most $\gamma$
each, with arbitrarily small slack.  Thus
\begin{equation}\label{eq:W-cb-near-normalizer}
 \dist\bigl(W(t)^{(k)}xW(t)^{(k)*},M_k(M)_1\bigr)\leq2\gamma.
\end{equation}
Here $T^{(k)}=I_{\C^k}\otimes T$.  Applying the same two approximations
before conjugating by $U_M(t)$ gives the corresponding assertion for
$W(t)^*xW(t)$.  Consequently
\begin{equation}\label{eq:W-range-gap}
 \dcb\bigl(W(t)MW(t)^*,M\bigr)\leq2\gamma
 \qquad(t\in\R).
\end{equation}

\paragraph{The frequency nest and its signs.}
Use the notation $\mathcal K,D,P_a$ introduced above.  Let $C_s$ denote
multiplication by $e^{ist}$, and let $\mathcal U_P$ be multiplication by
$U_P(t)$.  Spectral calculus gives
\[
 \mathcal U_P^*D\mathcal U_P=D+L_P,\qquad C_s^*DC_s=D+s.
\]
The sum on the right is the selfadjoint operator obtained from the
commuting spectral measures on the two tensor factors; these identities
therefore include their domains.  They can also be checked in the
spectral representation of $L_P$, where $\mathcal U_P$ is a scalar
frequency translation on each fiber.
Put $\mathcal W=\mathcal U_N\mathcal U_M^*$.  Directly moving the spectral
projections past these unitaries gives
\begin{equation}\label{eq:shifted-nest-block}
 (1-P_a)C_1\mathcal W P_a
 =C_1\mathcal U_N
   1_{(-\infty,a-1]}(D+L_N)
   1_{(a,\infty)}(D+L_M)\mathcal U_M^*.
\end{equation}
After Fourier transformation the product in the middle has fiber
\[
 1_{(-\infty,a-1-\xi]}(L_N)
 1_{(a-\xi,\infty)}(L_M).
\]
Its two thresholds have separation one, with midpoint
$a-\tfrac12-\xi$.  The cut estimate of \Cref{cor:modular-cuts} bounds its
norm by $C\gamma$, uniformly in $a$ and $\xi$.  The same argument with
$M,N$ interchanged applies to $C_1\mathcal W^*$.  In particular
\begin{equation}\label{eq:two-nest-errors}
 \sup_a\norm{(1-P_a)C_1\mathcal W P_a}\leq C\gamma,
 \qquad
 \sup_a\norm{(1-P_a)C_1\mathcal W^*P_a}\leq C\gamma.
\end{equation}

Adjoin $0,I$ to the complete nest $\{P_a:a\in\R\}$.
Arveson's distance formula for its nest algebra
\cite[Theorem~1.1, p.~212]{Arveson1975} says that the distance of an
operator $T$ to the operators preserving every $P_a\mathcal K$ equals
$\sup_a\norm{(1-P_a)TP_a}$.  Choose preserving operators $T_B,T_S$ with
\[
 \norm{T_B-C_1\mathcal W}\leq\delta,
 \qquad\norm{T_S-C_1\mathcal W^*}\leq\delta,
 \qquad \delta=C\gamma,
\]
enlarging $C$ to allow slack in taking the infimum.

\paragraph{Character averaging produces multipliers.}
The averaging below parallels Arveson's construction of an expectation
onto multiplication operators that preserves the frequency nest algebra
on $L^2(\mathbb T)$ \cite[Proposition~5.1, pp.~228--229]{Arveson1975}.
Use the invariant mean on the discrete group $\R$ from
\Cref{lem:invariant-mean}. For $T\in\B(\mathcal K)$ put
\[
 \mathcal E(T)=m_s C_sTC_s^*.
\]
This UCP map has range in the commutant of the scalar characters. Since
$C_s^*P_aC_s=P_{a-s}$, conjugation by $C_s$ preserves the entire nest
algebra.  The coefficient average preserves it as well: each identity
$(1-P_a)C_sTC_s^*P_a=0$ remains zero after taking the mean.
Both target multipliers $C_1\mathcal W$ and $C_1\mathcal W^*$ commute
with every $C_s$, so averaging retains both approximation bounds.

The scalar characters generate the scalar multiplication algebra
$L^\infty(\R)$.  Indeed, an $L^1$ function annihilating all characters
has zero Fourier transform and hence is zero.  Their commutant on
$L^2(\R,H)$ is therefore the algebra of decomposable operators.  One
may see the latter assertion by taking matrix coefficients relative
to a countable orthonormal basis of $H$: each coefficient commutes with
scalar multiplication and is itself a scalar multiplier; the operator
norm bound on a countable dense set of finite vectors gives a common
almost-everywhere bounded operator-valued field.  Thus
$\mathcal E(T_B)$ and $\mathcal E(T_S)$ are multipliers by fields
$B(s),S(s)$ satisfying
\begin{equation}\label{eq:B-S-boundary-errors}
 \mathop{\rm ess\,sup}_s\norm{B(s)-e^{is}W(s)}\leq\delta,
 \qquad
 \mathop{\rm ess\,sup}_s\norm{S(s)-e^{is}W(s)^*}\leq\delta.
\end{equation}
They preserve $P_0\mathcal K$, so \Cref{lem:analytic-multiplier} applies.
Write their extensions as $B(z),S(z)$.

\paragraph{Analytic products retain the approximate algebra relations.}
The boundary products satisfy, in both orders,
\[
 \norm{B(s)S(s)-e^{2is}I},\ 
 \norm{S(s)B(s)-e^{2is}I}\leq2\delta+\delta^2
 \quad\hbox{almost everywhere}.
\]
By the product identity and the positive Poisson kernel,
\begin{equation}\label{eq:analytic-inverse-errors}
 \norm{B(z)S(z)-e^{2iz}I},\ 
 \norm{S(z)B(z)-e^{2iz}I}\leq2\delta+\delta^2.
\end{equation}
Similarly, for every $k$ and $x\in M_k(M)_1$,
\begin{equation}\label{eq:analytic-algebra-errors}
 \dist\bigl(B(z)^{(k)}xS(z)^{(k)},M_k(M)\bigr)\leq\beta,
 \quad
 \dist\bigl(S(z)^{(k)}xB(z)^{(k)},M_k(M)\bigr)\leq\beta,
 \qquad \beta=2\gamma+2\delta+\delta^2.
\end{equation}
Here is a justification that avoids choosing measurable nearest
points in $M_k(M)$.  On the boundary, the inequalities follow from
\eqref{eq:W-cb-near-normalizer} and
\eqref{eq:B-S-boundary-errors}.  Test against any normal functional
annihilating $M_k(M)$.  The product identity identifies its value in
the upper half-plane with the Poisson average of its boundary values,
so the same bound holds there.  Since $M_k(M)$ is weak-* closed, its
norm distance is the supremum over these annihilating normal
functionals of norm at most one.  This last distance formula follows
from Hahn--Banach applied to the preannihilator in the trace class.

For $z=t+ir$ with $0<r\leq\gamma$ define
\[
 A_r(t)=e^{-iz}B(z),\qquad C_r(t)=e^{-iz}S(z).
\]
Because $|e^{-iz}|=e^r$, equations
\eqref{eq:analytic-inverse-errors}--\eqref{eq:analytic-algebra-errors}
give, uniformly in $r,t,k$,
\begin{align}
 \norm{A_r(t)},\ \norm{C_r(t)}&\leq e^\gamma(1+\delta)
                                      \leq1+C\gamma,\label{eq:AC-norms}\\
 \norm{A_r(t)C_r(t)-I},\ \norm{C_r(t)A_r(t)-I}
                                      &\leq C\gamma,\label{eq:AC-inverses}\\
 \dist\bigl(A_r(t)^{(k)}xC_r(t)^{(k)},M_k(M)\bigr),\quad
 \dist\bigl(C_r(t)^{(k)}xA_r(t)^{(k)},M_k(M)\bigr)
                                      &\leq C\gamma.
                                      \label{eq:AC-algebras}
\end{align}
The last line is for $x\in M_k(M)_1$.

\paragraph{Polar smoothing gives a norm-continuous path.}
For small enough $\gamma$, both products in \eqref{eq:AC-inverses} are
invertible.  They provide a right and a left inverse for $A_r(t)$,
which agree, and
\[
 \norm{A_r(t)^{-1}}\leq\frac{1+C\gamma}{1-C\gamma}
                                      \leq1+C\gamma.
\]
Together with \eqref{eq:AC-norms}, this places the spectrum of
$|A_r(t)|$ in $[1-C\gamma,1+C\gamma]$.  Its polar unitary
\[
 V_r(t)=A_r(t)\bigl(A_r(t)^*A_r(t)\bigr)^{-1/2}
\]
therefore satisfies
\begin{equation}\label{eq:polar-AC-errors}
 \norm{V_r(t)-A_r(t)}\leq C\gamma,
 \qquad \norm{V_r(t)^*-C_r(t)}\leq C\gamma.
\end{equation}
For the second bound, first use
$\norm{C_r(t)-A_r(t)^{-1}}
 \leq\norm{A_r(t)^{-1}}\norm{A_r(t)C_r(t)-I}$,
and then compare $A_r(t)^{-1}$ with $V_r(t)^*$.
It follows from \eqref{eq:AC-algebras} that $V_r(t)$ approximately
normalizes $M$ in both directions, with error $C\gamma$ at every
matrix level.

For fixed $r>0$, the fields $A_r(t)$ are norm continuous by
holomorphy, and their inverses are uniformly bounded.  Continuous
functional calculus makes $V_r(t)$ norm continuous.  Put
\[
 E_r(t)=V_r(t)V_r(0)^*.
\]
This is a norm-continuous unitary path with $E_r(0)=I$ and
\begin{equation}\label{eq:Er-normalizer-error}
 \sup_{x\in M_1}\dist(E_r(t)xE_r(t)^*,M)\leq C\gamma.
\end{equation}
To check composition here, approximate $V_r(0)^*xV_r(0)$ by an
$m\in M$ with $\norm{m}\leq1+C\gamma$, and apply the corresponding
estimate for $V_r(t)$ to $m$; the sum of the two errors is still
$C\gamma$.  Apply \Cref{lem:local-normalizer-path} to the portion of
this path joining $0$ to $t$, in either time direction.  For each
$r,t$ it gives $a_r(t)\in\U(M)$ and $b_r'(t)\in\U(M')$ such that
\begin{equation}\label{eq:Er-product-approximation}
 \norm{E_r(t)-a_r(t)b_r'(t)}\leq C\gamma,
\end{equation}
where a further slack of $\gamma$ allows for an unattained infimum.
No continuous choice of these factors is needed.

\paragraph{The boundary comparison holds at every fixed time.}
Set $F_0(s)=e^{is}W(s)$, a bounded strongly-* continuous field, and
\[
 D_r(t)=e^{-i(t+ir)}\int_\R
             \frac{rF_0(s)}{\pi((t-s)^2+r^2)}\,ds.
\]
The Poisson formula and \eqref{eq:B-S-boundary-errors} give
$\norm{A_r(t)-D_r(t)}\leq e^r\delta$.  Thus
\begin{equation}\label{eq:V-D-error}
 \norm{V_r(t)-D_r(t)}\leq C\gamma,
 \qquad \norm{D_r(t)}\leq e^r.
\end{equation}
At every fixed $t$, the Poisson approximate identity gives
$D_r(t)\to W(t)$ strongly-* as $r\downarrow0$.
For clarity, on a fixed vector split the integral into
$|s-t|<h$ and its complement.  Strong continuity controls the first
piece, while the Poisson mass of the complement tends to zero.
Apply the same argument to the adjoint field; the scalar prefactor
converges to $e^{-it}$.  This proves both strong limits everywhere,
not just almost everywhere.

Consequently
\[
 Q_r(t)=D_r(t)D_r(0)^*\longrightarrow W(t)
 \quad\hbox{strongly-* at every fixed }t,
\]
since $W(0)=I$ and all factors are uniformly bounded.  By
\eqref{eq:V-D-error},
\begin{equation}\label{eq:Er-Qr-error}
 \norm{E_r(t)-Q_r(t)}\leq C\gamma.
\end{equation}
These estimates compare $E_r(t)$ with a convergent family; they do
not assert convergence of $E_r(t)$ itself.

\paragraph{A UCP limit recovers an actual normalizer.}
Fix $t$.  For a sequence $r\downarrow0$ define the UCP maps
\[
 \Phi_r:M\longrightarrow M,\qquad
 \Phi_r(x)=a_r(t)x a_r(t)^*.
\]
Since $b_r'(t)$ commutes with $M$, this is the restriction of
$\Ad(a_r(t)b_r'(t))$.  Equations
\eqref{eq:Er-product-approximation} and \eqref{eq:Er-Qr-error}
therefore imply, at every matrix level,
\begin{equation}\label{eq:Phi-Qr-cb}
 \norm{\Phi_r^{(k)}(x)-Q_r(t)^{(k)}xQ_r(t)^{(k)*}}
        \leq C\gamma\norm{x},\qquad x\in M_k(M).
\end{equation}
Indeed, for any bounded operators $T,R$ the cb norm of the difference
of the maps $x\mapsto TxT^*$ and $x\mapsto RxR^*$ is at most
$(\norm{T}+\norm{R})\norm{T-R}$.

Compactness of the product of the weak-operator compact balls
$\{y\in M:\norm{y}\leq\norm{x}\}$, indexed by $x\in M$, supplies
a subnet for which $\Phi_r(x)$ converges weakly for every $x$.
Denote its limit by $\Phi(x)$.  Linearity, unitality and positivity
at every matrix level pass to this limit, so $\Phi:M\to M$ is UCP.
No normality of $\Phi$ is needed.  Bounded strong-* continuity of
multiplication gives
$Q_r(t)xQ_r(t)^*\to W(t)xW(t)^*$ strongly-* for each fixed $x$,
and likewise at every fixed matrix level.  Weak-operator lower
semicontinuity of norm in \eqref{eq:Phi-Qr-cb} now yields
\[
 \norm{\Phi-\rho_t}_{\cb}\leq C\gamma,
 \qquad \rho_t(x)=W(t)xW(t)^*.
\]
The representation $\rho_t$ is faithful and normal, and its range
has two-sided cb gap at most $2\gamma$ from $M$ by
\eqref{eq:W-range-gap}.  In particular the reverse ordinary gap
required by \Cref{prop:ucp} is small.  That proposition supplies a
unitary $q_t$ with
\[
 \norm{q_t-I}\leq C\gamma,\qquad
 q_t\rho_t(M)q_t^*=M.
\]
Then $n_t=q_tW(t)$ is the required unitary normalizer and
$\norm{n_t-W(t)}=\norm{q_t-I}\leq C\gamma$.
All estimates and smallness thresholds used only universal constants,
which proves the uniform assertion.
\end{proof}


\subsection{Correction within the normalizer group}

The preceding construction produces pointwise normalizers close to
a modular group. Their selections need not be measurable or continuous.
We correct their multiplication law using an invariant mean on the
discrete group $\R$. The final comparison with the original modular
group will then supply strong continuity.

The quadratic correction follows the averaging method for approximate
unitary representations \cite{Kazhdan1982,BurgerOzawaThom2013}.
Here the extra requirement is that every corrected unitary still
normalizes the prescribed algebra. The local factorization proved next
ensures that the averaging iteration preserves this requirement.

\begin{lemma}[Small exact normalizers]
\label{lem:small-exact-normalizer}
There are absolute $r_0,C_0>0$ such that, if $P$ is a separably
acting type $\III$ factor, $v\in\mathcal N_P$, and
$\norm{v-I}<r_0$, then
\[
 v=ab',\qquad a\in\U(P),\quad b'\in\U(P'),
 \qquad \norm{a-I}+\norm{b'-I}\le C_0\norm{v-I}.
\]
\end{lemma}

\begin{proof}
Write $v=e^{ih}$ by the principal logarithm, with $h=h^*$ and
$\norm{h}\le2\norm{v-I}$. For $r_0$ small, the norm-convergent
logarithm series on $\B(\B(H))$ gives
\[
 \log(\Ad_v)=i\ad_h.
\]
Indeed $\Ad_v=\exp(i\ad_h)$ and $\norm{\ad_h}$ is within the
radius where logarithm and exponential are inverse analytic
functions. Each term of the logarithm series preserves $P$,
since $vPv^*=P$. Hence $i\ad_h$ restricts to a bounded
$*$-derivation of $P$.

The factor $P$ is a simple unital $C^*$-algebra, by the same
projection argument as in \Cref{lem:local-normalizer-linear}.
Sakai's theorem therefore gives a self-adjoint $a_0\in P$ such
that $h-a_0\in P'$. A scalar may be subtracted from $a_0$ and
added to this commutant element. Choose the scalar at the midpoint
of the spectrum of $a_0$. We claim that the resulting $a_0$ has
$\norm{a_0}\le\norm{h}$.

For any self-adjoint $a\in P$, its spectral diameter is at most
$\norm{\ad_a|_P}$. In fact, choose nonzero spectral projections
near the two spectral endpoints and a norm-one partial isometry
between nonzero subprojections of them. Compressing $[a,w]$ by
$w^*$ shows that its norm is at least the endpoint separation
minus the two spectral tolerances. Let those tolerances tend to
zero. Apply this to $a_0$, using
$\norm{\ad_{a_0}|_P}=\norm{\ad_h|_P}\le2\norm{h}$.
After the midpoint shift, $\norm{a_0}\le\norm{h}$ and
$b_0'=h-a_0$ has norm at most $2\norm{h}$.
They commute, so $v=e^{ia_0}e^{ib_0'}$; these are the required
factors, with total displacement at most $3\norm{h}$.
\end{proof}

In the next proof all operator means use \Cref{lem:invariant-mean}
for the discrete group $\R$. Thus arbitrarily selected factors may
be averaged separately in $P$ and $P'$, without a measurability assumption.

\begin{lemma}[Correction of normalizer representations]
\label{lem:normalizer-representation-correction}
There are absolute $e_1,C_1>0$ such that the following holds.
Let $P\subset\B(H)$ be a separably acting type $\III$ factor.
For any family $v_t\in\mathcal N_P$, $t\in\R$, with
\[
 e:=\sup_{t,s\in\R}\norm{v_tv_s-v_{t+s}}<e_1,
\]
there is an exact unitary representation $z:\R\to\mathcal N_P$
of the discrete group with
\[
 \sup_t\norm{z_t-v_t}\le C_1e.
\]
\end{lemma}

\begin{proof}
Let $m_s$ denote the invariant operator mean just described. Put
\[
 F_t(s)=v_{t+s}v_s^*,\qquad A_t=m_sF_t(s).
\]
Every $F_t(s)$ is a unitary normalizer and
$\norm{F_t(s)-v_t}\le e$. Hence $\norm{A_t-v_t}\le e$.
Translation invariance and the exact factorization
\[
 F_{t+r}(s)=F_t(r+s)F_r(s)
\]
give
\begin{align*}
 A_{t+r}-A_tA_r
 &=m_s\bigl[(F_t(r+s)-A_t)(F_r(s)-A_r)\bigr],\\
 \norm{A_{t+r}-A_tA_r}&\le4e^2.
\end{align*}
The centered terms have norm at most $2e$, and their means vanish;
expanding verifies the first equality with the written operator
order. Likewise
\[
 I-A_t^*A_t
 =m_s(F_t(s)-A_t)^*(F_t(s)-A_t),
\]
so its norm is at most $4e^2$. Taking polar parts in $\B(H)$
alone would not ensure membership in $\mathcal N_P$. We instead
use the local product structure of exact normalizers.

By \Cref{lem:small-exact-normalizer}, choose, for every $s,t$,
\[
 F_t(s)v_t^*=a_{t,s}b_{t,s}',
 \quad a_{t,s}\in\U(P),\quad b_{t,s}'\in\U(P'),
 \quad
 \norm{a_{t,s}-I}+\norm{b_{t,s}'-I}\le C_0e.
\]
All these choices can be arbitrary. Their averages
$a_t=m_sa_{t,s}\in P$ and $b_t'=m_sb_{t,s}'\in P'$ satisfy
\[
 \norm{m_s(a_{t,s}b_{t,s}')-a_tb_t'}\le4C_0^2e^2.
\]
This follows by expanding the centered product; its two factors
have norm at most $2C_0e$. Unitarity of the individual factors
also gives
\[
 \norm{I-a_t^*a_t}\le4C_0^2e^2,
 \qquad \norm{I-(b_t')^*b_t'}\le4C_0^2e^2.
\]
For small $e$ the averages are invertible, since they are within
$C_0e$ of $I$. Their polar unitaries $\widehat a_t\in P$ and
$\widehat b_t'\in P'$ differ from them by at most $Ce^2$, with
an absolute $C$: use $\norm{|a_t|-I}\le
\norm{a_t^*a_t-I}/(1+\sqrt{1-4C_0^2e^2})$, and similarly on
the other side. Therefore
\[
 w_t=\widehat a_t\widehat b_t'v_t\in\mathcal N_P,
 \qquad \norm{w_t-A_t}\le C_2e^2.
\]
It follows, uniformly in $t,r$, that
\[
 \norm{w_t-v_t}\le e+C_2e^2,
 \qquad
 \norm{w_tw_r-w_{t+r}}\le(4+3C_2)e^2.
\]
For the second inequality use $\norm{A_t}\le1$ and
$\norm{w_t}=1$ to compare products, then the preceding bound
on the defect of $A$.

Iterate this operation, always with the same absolute constants.
Choose $e_1$ so small that every new defect is at most half the
preceding one and every displacement is at most twice the preceding
defect. The total displacement is then at most $4e$.
The families converge uniformly in $t$ in operator norm to
unitaries $z_t$, satisfying $z_tz_r=z_{t+r}$.
The group $\mathcal N_P$ is norm closed: if unitaries $u_j\to u$
normalize $P$, norm closure of $P$ gives $uPu^*\subset P$ and,
by applying the same argument to $u_j^*$, the reverse inclusion.
Thus $z_t\in\mathcal N_P$ for all $t$. The exact multiplication
law implies $z_0=I$ and $z_{-t}=z_t^*$, as required.
\end{proof}

\begin{proof}[Proof of \Cref{prop:modular-alignment}]
Put $U_t=\Delta_N^{it}$ and
$W(t)=\Delta_N^{it}\Delta_M^{-it}$.
\Cref{lem:modular-near-normalizers} supplies $n_t\in\mathcal N_M$
with $\sup_t\norm{n_t-W(t)}\le C\gamma$. Set
$v_t=n_t\Delta_M^{it}$, so $v_t\in\mathcal N_M$ and
$\sup_t\norm{v_t-U_t}\le C\gamma$. Since $U$ is an exact
representation, $\sup_{t,s}\norm{v_tv_s-v_{t+s}}\le3C\gamma$.
For sufficiently small universal $\gamma$, apply
\Cref{lem:normalizer-representation-correction} to obtain an
exact representation $z_t\in\mathcal N_M$ with
\[
 \sup_t\norm{z_t-U_t}\le C'\gamma.
\]

Take the invariant mean
$T=m_tz_tU_t^*$. It is within $C'\gamma$ of $I$ and hence
invertible for small $\gamma$. For every $s$, translation
invariance and the two exact representation laws imply
\[
 z_sTU_s^*=m_tz_{s+t}U_{s+t}^*=T.
\]
Its polar unitary $w=T|T|^{-1}$ therefore also satisfies
$z_sw=wU_s$. Moreover
$\norm{w-I}\le2\norm{T-I}\le2C'\gamma$.
In particular $z_s=wU_sw^*$; this proves strong continuity of $z$
without any regularity assumption on the selected normalizers.

Set $q=w^*$. Since $z_s$ normalizes $M$, the intertwining identity
gives $U_s(qMq^*)U_s^*=qMq^*$ for every $s$.
At each matrix level conjugation by $q^{(k)}$ moves contractions
by at most $2\norm{q-I}$. Consequently
\[
 \dcb(qMq^*,N)\le\gamma+2\norm{q-I}\le C_{\rm mod}\gamma.
\]
Enlarging $C_{\rm mod}$ and reducing $c_{\rm mod}$ to meet the
finitely many universal smallness conditions proves all assertions.
\end{proof}
